\documentclass[conference]{IEEEtran}
\usepackage[utf8]{inputenc}
\usepackage[T1]{fontenc}
\usepackage{cite}
\usepackage{amsmath,amssymb}
\usepackage{graphicx}
\usepackage{booktabs}
\usepackage{url}

\title{Verifier-Guided Prompting for Intent-to-Configuration Translation: A Preregistered NetOps Study}

\author{
\IEEEauthorblockN{Autonomous AI Researcher}
\IEEEauthorblockA{CNSM 2026 Agentic AI Researcher Track}
}

\begin{document}

\maketitle

\begin{abstract}
We preregistered and executed a paired confirmatory trial to test whether constrained prompting plus a deterministic verifier-guided generate\ensuremath{\rightarrow}verify\ensuremath{\rightarrow}repair loop (one allowed repair) increases deterministic-verifier semantic-correctness when translating operator natural-language intents into low-level device configurations. Design: N=150 synthetic intents sampled from a deterministic generator, shared\_initial\_candidate generation semantics, model = gpt-5-mini, deterministic validator = deterministic\_netops\_validator\_v5, one initial generation per intent shared across baseline and guarded conditions, guarded pipeline permitted <=1 repair call. Results (artifact-grounded): baseline = 149/150 (99.333\%); guarded = 150/150 (100.0\%); paired absolute difference = 0.006667 (\ensuremath{\approx}0.67 percentage points); exact McNemar two-sided p = 1.0; 95\% bootstrap percentile CI = [0.00, 0.02] (10,000 resamples, seed = 1729). Provider accounting: completed episodes = 300; model\_calls\_used = 151; repair-stage invocations = 1. Interpretation: on this preregistered synthetic corpus and under shared\_initial\_candidate semantics the guarded workflow produced a numerically negligible and statistically non-significant improvement over a near-ceiling baseline. We emphasize the methodological lesson that preregistered NetOps trials require baseline-calibration pilots to avoid ceiling effects that invalidate preregistered power assumptions. All execution and analysis artifacts are archived in the supplied execution bundle referenced by the execution manifest.
\end{abstract}

\section{Introduction}
Natural-language intent interfaces aim to reduce operator burden by letting humans express desired network changes as free text that is automatically translated into executable device configurations. Prior work documents that single-pass LLM outputs often contain syntactic errors, omissions, or semantic violations that can yield unsafe or unavailable states if applied directly [1][2][3]. A commonly proposed defense is a generate\ensuremath{\rightarrow}verify\ensuremath{\rightarrow}repair architecture in which an LLM produces a candidate, a deterministic verifier checks intent-level invariants or formalized constraints, and an automated repair (or human gate) corrects violations before execution [4][5]. Security studies of tool-augmented LLM agents further motivate deterministic gating because unchecked textual claims can become harmful actions in multi-tool workflows [6].

Research question. After an autonomous literature review and preregistration we posed a confined confirmatory question: does constrained prompting (structured templates with explicit semantic slots) combined with deterministic verifier-guided iterative feedback materially increase verifier-level semantic correctness of NL\ensuremath{\rightarrow}configuration translation versus an unconstrained baseline under a paired design? The trial (study\_cand\_2026\_b2\_v1) was preregistered and frozen before any model outputs were observed.

Contributions. (1) A fully preregistered, artifactized paired confirmatory trial measuring deterministic-verifier pass/fail on a programmatically generated synthetic corpus (N=150 paired intents). (2) Artifact-backed outcomes showing that, under shared\_initial\_candidate semantics and a one-repair budget, the guarded pipeline raised verifier pass rate from 149/150 to 150/150 (+0.67 pp), a numerically negligible and statistically non-significant improvement (exact McNemar p = 1.0; 95\% bootstrap CI = [0.00, 0.02]). (3) A methodological recommendation that preregistered NetOps trials include baseline-calibration pilots and explicitly specify generation semantics (shared\_initial\_candidate vs independent sampling) because semantics materially change the estimand and interpretation.

\section{Related Work}
Related literature spans intent-translation evaluations, verifier-guided repair systems, formal verifier tooling, and agentic/tool-augmented safety analyses. First, multiple recent studies evaluating NL\ensuremath{\rightarrow}configuration or intent-to-configuration pipelines report that translating underspecified operator language into verifier-sound artifacts is challenging in practice. Intent-driven translation work documents that free-text prompts frequently produce omissions or incorrect sequencing that violate required invariants, and published benchmark and evaluation studies report modest semantic-correctness rates on heterogeneous task banks [1]. Those evaluations emphasize that the mapping from natural intent to an exact required command sequence depends critically on prompt design, model competence, and the chosen correctness oracle.

Second, a growing body of work integrates deterministic verification into generation workflows and reports improved practical correctness, but gains depend strongly on baseline competence, verifier expressiveness, and repair budget. Systems that use policy-guided verifier feedback or verifier-in-the-loop fine-tuning (TerraFormer and related IaC efforts) demonstrate that verifier signals can be used to rerank or fine-tune candidate generations and thereby reduce syntactic and semantic errors for infrastructure-as-code artifacts [2]. Independent benchmark efforts focused specifically on LLM-driven configuration repair likewise show that generate\ensuremath{\rightarrow}verify\ensuremath{\rightarrow}repair pipelines can improve downstream validity, but reported effect sizes vary across datasets and experimental semantics (single-shot vs multi-sample generation, number of allowed repairs) and across whether the verifier feeds gradient-based updates or only deterministic textual feedback to a repair prompt [3]. These method papers collectively highlight two operational knobs that we treat explicitly in our design: the repair budget (how many additional generations/repairs are permitted) and the sampling semantics (whether different conditions draw independent first-pass samples or share an initial candidate). Both knobs materially affect the estimand and therefore the interpretable marginal benefit of verifier-guided repair.

Third, deterministic formal-verification tools provide concrete oracles for many network-level invariants but carry limitations that affect their role as evaluation endpoints. Fast symbolic/network-verification engines (e.g., NetKAT/KATch family) and related packet-reachability checkers can deterministically validate reachability, isolation, and many high-level policy constraints and have been proposed as suitability oracles for LLM-generated configurations [5]. Those tools are valuable as audit or gating layers because they return compact diagnostic counters and counterexamples that can drive automated repairs. At the same time, tool limitations---supported protocol models, topology scale, and abstraction mismatch with vendor-specific CLI sequences---mean that verifier coverage and mapping correctness become additional sources of error. Prior systems therefore combine careful mapping/unit-tests with verifier-coverage checks before committing to large-batch automated evaluation [5].

Fourth, security and systems analyses of agentic LLM pipelines warn that unguarded multi-tool agents can convert false textual claims into harmful actions and are susceptible to prompt-injection and claim-to-action exploit patterns [4]. Empirical attack studies show that deployed-style tool orchestration layers can be manipulated when the agent's decision pathways are not auditable or when verifier gating is absent; these findings motivate designs that make verifier verdicts explicit, auditable, and conservative in the presence of ambiguous intents.

How this study situates. Our trial operationalizes a narrowly constrained guarded architecture and emphasizes preregistration, paired inference, and deterministic-verifier endpoints. It differs from many prior reports in three design choices that prior work identifies as consequential: (a) we enforced shared\_initial\_candidate execution semantics so that baseline and guarded comparisons are computed on the identical sampled candidate (this isolates the marginal value of a repair applied to a fixed first-pass output); (b) we constrained the repair budget to at most one repair per intent, consistent with low-latency operational budgets considered in some IaC workflows; and (c) we used a deterministic verifier as the strict binary oracle for semantic correctness. These choices intentionally reduce the maximal observable marginal benefit of guarded repair under our estimand, explaining how previously reported larger improvements can coexist with our bounded result when different sampling or repair regimes are evaluated [2][3][5][6].

\section{Methodology}
Preregistration and confirmatory design. The study was preregistered as study\_cand\_2026\_b2\_v1 with a frozen analysis plan executed before any model outputs were observed. The primary estimand was the paired\_success\_rate\_difference\_guarded\_minus\_baseline evaluated with an exact two-sided McNemar test on deterministic-verifier pass/fail outcomes. The confirmatory sample specified N=150 distinct synthetic intents drawn deterministically from a preregistered generator, deterministic inclusion/exclusion rules (a rule-based ambiguity detector and verifier-coverage checks), a deterministic retry policy for provider timeouts (one retry), and a power target to detect an absolute +15 percentage-point improvement (\ensuremath{\approx}80\% power) under conservative discordant-pair assumptions. All elements of the preregistration (estimands, inclusion/exclusion rules, multiplicity plan, and analysis scripts) were frozen and archived prior to model calls.

Execution semantics: shared\_initial\_candidate, call budgets, and adapter enforcement. A central design decision was the execution-semantic enforcement of shared\_initial\_candidate: for each intent the execution adapter produced exactly one initial model generation and that identical sampled candidate was used in the paired evaluation of both baseline and guarded conditions. The guarded pipeline was permitted at most one additional repair call only if the deterministic verifier failed that initial candidate. Operationally, this means the baseline rate measures the validity of the single initial sampled candidate; the guarded vs baseline contrast therefore isolates the marginal effect of applying at most one repair to that identical candidate rather than measuring performance averaged over independent first-pass draws. The adapter enforced these constraints at runtime and recorded stage-level provider accounting; archived provider audits report stage\_counts = \{shared\_initial\_generation: 150, repair: 1\} and model\_calls\_used = 151, consistent with the enforced one-initial plus at most-one-repair policy.

Model, sampling, and deterministic validator. Declared model: gpt-5-mini (provider record: openai\_responses). The execution metadata include an explicit record that model\_configuration.json declares declared\_model\_version = "gpt-5-mini" and temperature = null (unset). We emphasize the literal meaning: temperature = null in the recorded configuration is not equivalent to setting temperature = 0 (deterministic decoding). Because no numeric temperature value or fixed random-seed was enforced in the preregistration or execution, model outputs may be nondeterministic across independent API sessions. This nondeterminism was anticipated in the preregistration and is reported explicitly here because sampling semantics materially affect estimand interpretation. Deterministic verifier: deterministic\_netops\_validator\_v5 (validator\_version = 5.0) applied a full\_intent\_constraint\_validation rule and returned for each candidate a boolean pass/fail verdict plus normalized final-configuration text and per-episode diagnostics (parsed\_command\_count, required\_command\_count, observed\_touched\_field\_count, violation\_codes). These deterministic verdicts constituted the binary endpoint for the confirmatory test.

Task generator, corpus, prechecks, and mapping verification. Confirmatory tasks were produced deterministically by the preregistered generator (deterministic\_netops\_task\_generator\_v5) and were stratified across four intent families: reachability/isolation, ACL-like access changes, VLAN/MTU sequence changes, and uplink switchover patterns. Topologies were constrained to verifier-capable small networks (4--32 nodes) to ensure deterministic oracle coverage. Prior to batch execution we ran the preregistered mapping and verifier-coverage unit-tests: parser transformations that map LLM textual outputs to verifier inputs were unit-tested and passed. The preregistration specified deterministic exclusion rules (ambiguous intents routed to exploratory corpus; parse failures or unsupported verifier features excluded with deterministic reason codes); no confirmatory exclusions occurred in the archived run. Representative generator seeds, task\_payloads, and expected\_final\_states are archived to support independent calibration or pilot reuse.

Scoring rules, primary analysis, and confidence procedures. The primary outcome was deterministic verifier pass (binary). Scoring used the archival scorer identifier deterministic\_netops\_validator\_v5 and its full\_intent\_constraint\_validation rule; for each episode the scorer produced a normalized\_configuration string that was compared to the preregistered reference invariants for automated pass/fail determination. The confirmatory analysis used an exact two-sided McNemar test on the paired contingency (n\_11, n\_10, n\_01, n\_00) and computed a paired bootstrap percentile 95\% confidence interval with 10,000 resamples (bootstrap\_seed = 1729). The preregistered handling of incomplete pairs specified exclusion from the primary test (complete\_pair\_primary). Secondary contrasts (constrained vs baseline, guarded vs constrained) and multiplicity corrections were preregistered; under shared\_initial\_candidate enforcement these secondary contrasts were constrained by the adapter semantics and the available stage counts. All analysis code, contingency tables, and bootstrap parameters are archived and cited in the analysis artifacts.

Execution accounting, contamination controls, and reproducibility. Planned episodes were 300 (150 intents \ensuremath{\times} 2 conditions); completed\_episode\_count = 300 and complete\_pair\_count = 150 in the archived manifest. Provider-call accounting and contamination checks were recorded automatically: completed\_hosted\_model\_call\_event\_count = 151, failed\_hosted\_model\_call\_event\_count = 0, cache\_miss\_count = 151, and contamination\_summary.csv recorded zero flags. These runtime audits and per-call metadata (provider traces, shared\_initial\_response paths, and validator traces) were retained in the execution bundle to enable deterministic reconciliation of the paired counts and to support post hoc reproducibility checks. The preregistration also specified a contamination-detection checklist (session isolation, duplicate-response checks, mapping unit-tests) that the adapter enforced prior to batch execution; all checks passed and are documented in the artifact logs.

\section{Results}
Execution and provider audit. The run completed as planned: completed\_episode\_count = 300 and complete\_pair\_count = 150. Provider-call accounting: model\_calls\_used = 151; completed\_hosted\_model\_call\_event\_count = 151; failed\_hosted\_model\_call\_event\_count = 0; cache\_miss\_count = 151. Stage-level counts: shared\_initial\_generation = 150 and repair = 1. No contamination flags or pre-execution unit-test failures occurred.

Primary numerical outcomes and contingency counts. The paired contingency counts (guarded \ensuremath{\times} baseline) were: n\_11 = 149 (both-pass), n\_10 = 1 (guarded-only-pass), n\_01 = 0 (baseline-only-pass), n\_00 = 0 (both-fail). Per-condition success rates: baseline\_success\_count = 149/150 = 0.9933333333333333; guarded\_success\_count = 150/150 = 1.0. Estimated paired absolute difference (guarded - baseline) = 0.00666666666666671 (\ensuremath{\approx}0.67 percentage points). Exact McNemar two-sided p = 1.0. Paired bootstrap percentile 95\% CI = [0.00, 0.02] (10,000 resamples; bootstrap\_seed = 1729). These figures are reported verbatim from the archived analysis artifacts (analysis/results.json, analysis/paired\_contingency\_table.csv, analysis/condition\_summary.csv).

Repair diagnostics and revision iterations. Only one guarded repair invocation occurred (1/150), consistent with the near-ceiling initial success rate. That single repair converted one initially invalid candidate to a valid configuration under the deterministic verifier. Median revision iterations for guarded episodes = 0 (IQR = 0). Validator traces include validation\_before and validation\_after states, normalized\_configuration strings, parsed\_command\_count, required\_command\_count, violation\_codes, and a valid flag for each episode; representative traces for tasks 000001--000006 are archived and demonstrate varied difficulty patterns while confirming validator behavior.

Representative artifact-grounded example. Example pair (task-000001, generator\_seed = 1) intent: "Migrate edge1 to VLAN 30 and MTU 1600. For safety, shut edge1 down before changing either VLAN or MTU, then restore it to admin up. Do not modify any other interface." Shared-initial response and both condition responses returned the expected four-command sequence. Validator diagnostics for task-000001: parsed\_command\_count = 4, required\_command\_count = 4, observed\_touched\_field\_count = 3, and valid = true. The shared-initial candidate SHA-256 and normalized\_configuration strings are archived in the execution bundle as part of the scorer traces.

Sensitivity and missingness. The preregistered sensitivity analysis treating missing guarded outputs as failures was not invoked because failed\_hosted\_model\_call\_event\_count = 0 and unscorable episodes = 0. No exclusions or parse failures occurred and contamination\_summary.csv is empty.

Compact evidence tables. Table 1 (paired contingency) and Table 2 (execution audit summary) summarize the confirmatory counts and provider accounting; both tables are artifact-grounded and constructed from analysis/paired\_contingency\_table.csv and analysis/condition\_summary.csv.

Interpretation of numeric outcome. On this synthetic corpus and under shared\_initial\_candidate semantics the guarded workflow produced a numerically negligible absolute improvement (+0.67 pp) from an already near-ceiling baseline. The exact McNemar p = 1.0 and the bootstrap CI containing zero indicate no statistically supported confirmatory benefit under the preregistered estimand. Importantly, the observed baseline success rate (\ensuremath{\approx}99.33\%) violated the preregistered power assumption (target detectable effect = 15 pp), producing a ceiling effect that invalidated the originally planned sensitivity.

\section{Discussion}
Ceiling effect and implications for preregistration. The preregistered power calculation assumed a substantially lower baseline; observed baseline\_success\_rate = 0.9933 left almost no headroom for the guarded repair to show a large absolute improvement. This ceiling effect invalidated the preregistered detectable-effect assumptions (target = 15 percentage points) and reduced the trial's ability to demonstrate benefit. Concretely, with n\_11 = 149 and only a single discordant pair (n\_10 = 1), the paired analysis is dominated by an almost-complete concordance that leaves little information for estimating discordant behavior. We therefore recommend that preregistered NetOps trials include a short baseline-calibration pilot (for example, 20--40 tasks sampled from the planned generator and prompt style) to estimate baseline competence and either increase planned N or raise task difficulty before committing to confirmatory sampling. The archived deterministic task generator seeds and representative tasks permit such a pilot: one can sample the same generator with a small pilot, measure baseline\_success\_rate, and then choose confirmatory N or difficulty stratification so that the expected baseline leaves statistical headroom for the target absolute improvement.

Why the synthetic corpus produced a high baseline. The deterministic generator produced tasks with declared difficulty markers (medium, hard, very\_hard) and explicit required\_sequence constraints; inspection of representative traces (tasks 000001--000006) shows many tasks have fully specified required sequences and constrained action spaces (allowed\_fields, allowed\_touched\_fields) that reduce translation ambiguity. When intents are tightly specified and the verifier coverage matches the required invariants, even a single sampled candidate from a capable LLM often satisfies the verifier. In our run this interaction produced the near-ceiling baseline, illustrating that generator design choices (ambiguity, under-specification, and difficulty-sampling) materially influence the observable marginal benefit of any repair stage.

Semantics, sampling, and estimand trade-offs. The execution semantics (shared\_initial\_candidate) intentionally reduced per-intent sampling variance by using the same initial candidate across conditions; this made the guarded-baseline contrast a pure marginal-repair estimand. Alternative semantics---independent per-condition sampling or ensembles of initial draws---estimate different operational questions: independent sampling measures expected verifier pass rate when each pipeline draws independently from the model distribution, while multi-draw or multi-repair budgets measure how re-sampling or repeated repairs trade cost for improved correctness. Independent sampling or larger repair budgets typically increase the chance of observing discordant pairs (higher n\_10 + n\_01) and therefore can increase power to detect smaller absolute differences, but at the cost of additional model calls and latency. Practitioners should select sampling semantics that match deployment constraints (whether systems will resample or must act on a single first-pass candidate) and incorporate model-call cost and repair-latency budgets into experimental power planning.

Repair budgets, cost, and diminishing returns. The provider accounting from our run (model\_calls\_used = 151; repair-stage invocations = 1) quantifies the operational cost of the guarded design under the one-repair budget: repairs were rarely invoked when the baseline was near-ceiling, so marginal cost per observed success-gain was effectively high. In scenarios with lower baseline competence, repair budgets exhibit diminishing returns: initial repairs convert many failures, but additional repairs yield progressively fewer conversions. Quantifying diminishing returns requires multi-arm experiments that vary allowed repair counts and measure per-repair marginal gain versus model-call cost; our artifacts supply seeds and code to run such arms reproducibly.

Operational affordances beyond pass/fail. Deterministic guarded workflows provide benefits not captured by a single binary endpoint. The validator produces normalized\_configuration strings, parsed\_command\_count, required\_command\_count, and explicit violation\_codes that operators can audit; these artifacts accelerate triage and human review even when verifier-pass differences are numerically small. For example, representative validator traces in the archive show concrete violation\_code semantics (missing step, order-violation) that would guide automated or human repair policies. Thus, guarded pipelines can improve operational observability, reduce human triage time, and support auditable change records independent of marginal verifier-pass improvements on constrained corpora.

Threats to validity and residual risks. Internal validity: adapter-enforced call budgets and the provider audit mitigate cross-condition leakage, and mapping unit-tests reduced parse/translation failures. Construct validity: deterministic validator pass/fail is a proxy for semantic correctness and depends on mapping code and verifier coverage; although mapping unit-tests passed and no parse failures occurred, any uncovered verifier-coverage gap would bias semantic-correctness measurement. External validity: experiments used a single declared model (gpt-5-mini) and synthetic tasks on small topologies (4--32 nodes); results should not be extrapolated to larger, heterogeneous networks or other models. Security and adversarial robustness: this confirmatory trial did not evaluate adversarial prompt-injection; prior work indicates claim-to-action attacks remain an independent concern and defensive evaluation requires separate adversary-driven experiments. Conclusion validity: the mismatch between preregistered detectable-effect assumptions and the observed baseline underscores the importance of pilot calibration to preserve inferential validity.

Practical experimental prescriptions. Based on artifact-grounded outcomes we recommend: (1) a baseline-calibration pilot (20--40 tasks) to estimate baseline\_success\_rate and discordant-pair proportions prior to final N selection; (2) explicit preregistration of generation semantics and repair budgets; (3) multi-arm designs that vary repair budgets and sampling semantics to estimate cost--benefit frontiers; (4) inclusion of provider-call accounting in power and cost planning; and (5) staged scaling of verifier coverage and topology size with deterministic unit-tests for mapping code. The archived execution manifest and representative tasks provide concrete seeds and runnable code to implement these prescriptions reproducibly.

\section{Limitations}
\begin{itemize}
\item Single-model constraint: experiments used only gpt-5-mini; results do not generalize across models or sizes.
\item Synthetic corpus and scale: tasks were programmatically generated and limited to verifier-capable toy-to-small topologies (4--32 nodes); external validity to production WANs and heterogeneous vendor CLIs is untested.
\item Shared\_initial\_candidate semantics and execution contract limited repair budget to at most one guarded repair call per intent; different generation semantics or multiple-repair budgets may yield different outcomes.
\item Ceiling effect and preregistration mismatch: baseline correctness was far higher than the pilot assumptions used for power calculations (planned detectable effect = 15 percentage points), reducing the trial's ability to demonstrate benefit and illustrating sensitivity to baseline calibration.
\item Verifier coverage and specification translation: the study measures deterministic validator pass/fail on the preregistered invariants but does not claim safety guarantees beyond the deterministic validator's modeled properties.
\item Security and adversarial robustness: the experiment did not evaluate adversarial prompt-injection or adversary-driven intents; separate targeted studies are required to assess claim-to-action vulnerabilities in agentic NetOps workflows.
\end{itemize}

\section{Conclusion}
We executed a fully preregistered paired confirmatory trial comparing baseline free-text prompting with constrained-template prompting plus a single verifier-guided repair under shared\_initial\_candidate semantics. On a deterministic synthetic corpus with gpt-5-mini and deterministic\_netops\_validator\_v5, the guarded pipeline increased verifier pass rate from 149/150 to 150/150 (\ensuremath{\approx}0.67 pp), a numerically tiny and statistically non-significant difference (exact McNemar p = 1.0; 95\% bootstrap CI = [0.00, 0.02]). The principal scientific lesson is methodological: preregistered NetOps trials should include baseline calibration to avoid ceiling effects and preserve statistical power. Technically, our artifacts bound the marginal benefit of a one-repair guarded pipeline under shared\_initial\_candidate semantics; evaluating independent-sampling semantics, larger repair budgets, model diversity, and production-scale topologies remains important future work.

\section*{Disclosure Statement}
This study was executed autonomously after an autonomy lock under the archived master prompt (provenance/master\_prompt.txt, SHA-256 = 1872df1e\allowbreak{}1805d2d9\allowbreak{}6940456c\allowbreak{}a016bd66\allowbreak{}5d1d5196\allowbreak{}add77f5a\allowbreak{}cdf1582b\allowbreak{}b39b15ba). Preregistration identifier: study\_cand\_2026\_b2\_v1 (preregistration was frozen prior to observing outcomes and the preregistration record is archived). Declared model/tooling: declared model = gpt-5-mini (provider record: openai\_responses); execution adapter family = hosted\_netops\_gvr\_v1; deterministic validator = deterministic\_netops\_validator\_v5 (validator\_version = 5.0). Runtime sampling semantics (explicit): the archived model\_configuration.json records temperature = null (unset); because no numeric temperature or fixed random-seed was enforced, model outputs may be nondeterministic across independent API sessions. Autonomy and human role: a human Research Director selected the topic family and constrained the initial master prompt prior to the autonomy lock; after the lock the autonomous system performed literature review, study selection, preregistration, code generation, experiment execution, analysis, manuscript drafting, autonomous peer review, and revision. No human manually edited technical content, analyses, figures, captions, or references after the autonomy lock. Key execution facts reported in the paper (artifact-grounded): completed episodes = 300; complete paired tasks = 150; model\_calls\_used = 151; repair-stage invocations = 1. All runtime sampling metadata, provider-call traces, verifier inputs/verdicts, and analysis artifacts are archived in the execution bundle referenced by the execution manifest.

\begin{thebibliography}{99}
\bibitem{ref1} https://doi.org/10.1002/nem.2313
\bibitem{ref2} https://doi.org/10.1145/3786583.3786898
\bibitem{ref3} https://arxiv.org/abs/2604.22513
\bibitem{ref4} https://doi.org/10.1145/3605764.3623985
\bibitem{ref5} https://doi.org/10.1145/3656454
\bibitem{ref6} https://doi.org/10.1002/widm.70111
\end{thebibliography}

\end{document}
